\documentclass[11pt]{article}
\usepackage[utf8]{inputenc}
\usepackage[T1]{fontenc}
\usepackage[french]{babel}
\usepackage{amsmath,amssymb,amsthm}
\usepackage{geometry}
\usepackage{booktabs}
\usepackage{hyperref}
\geometry{margin=2.5cm}

\newcommand{\R}{\mathbb{R}}
\newcommand{\N}{\mathbb{N}}
\newcommand{\Z}{\mathbb{Z}}

\title{Instance numérique à 25 composants : maintenance d'un poste offshore\\
\large Génération de colonnes vs.\ contraintes en espérance}
\author{}
\date{}

\begin{document}
\maketitle

\section{Description de l'instance}

On reprend exactement le formalisme de \emph{Fluid approximation with column generation and an
application to offshore substations maintenance} --- le document qui décrit précisément le modèle codé
dans \texttt{MultiHorizonMaintenanceModel} --- pour construire une instance à 25 composants comparant
les deux méthodes de résolution : le LP maître restreint (RMP) résolu par génération de colonnes
(Méthode 1, Sections 3.3--3.4 du document de référence) et sa relaxation directe, dite à contraintes en
espérance (Méthode 2, Section 5).

\paragraph{Données de l'instance.}
\begin{itemize}
  \item $C = 25$ composants.
  \item Étapes (\emph{stages}) de $T=60$ jours (2 mois), comme dans le code existant. Avec un horizon
        de 20 ans et 6 étapes par an, $\bar{T} = 120$ étapes.
  \item Chaque composant a 3 états (1 = neuf, 2 = dégradé, 3 = en panne), MTBF $=40$ ans, identique pour
        tous les composants.
  \item Pas d'aléa météo : un seul scénario ($S=1$), le poste est toujours accessible et la production
        toujours à son maximum ($p_\tau \equiv 100$ chaque jour accessible). Une maintenance planifiée se
        déroule donc toujours exactement comme prévu, sans retard ni interruption.
  \item $R=3$ techniciens ; chaque opération de maintenance nécessite $r_c=1$ technicien, sans partage.
  \item Chaque composant $c$ est associé à une opération de maintenance qui le ramène à l'état 1, de durée
        $d_c \in \{1,\dots,10\}$ (entier tiré une fois pour toutes uniformément, donnée fixe de l'instance).
  \item Coût de remplacement $\mathrm{ccc}_c = 100$ pour tout $c$ (terme additif dans l'objectif, en
        plus de la pénalité de capacité).
  \item Quota annuel $K_0=10$ jours de maintenance gratuits, renouvelé tous les 6 étapes (1er janvier).
  \item Puisqu'il n'y a pas de météo, on peut toujours supposer que les plannings de maintenance
        commencent au jour 1 de l'étape (indifférence des jours).
\end{itemize}

\section{Coûts précalculés}

\paragraph{Transitions de dégradation.} Sans maintenance ($m=1$), $p_{c,x,x',1,t}$ est la matrice de
transition sur 60 jours obtenue, comme dans \texttt{build\_MC\_utils.jl}, en élevant à la puissance 60 la
matrice journalière absorbante linéaire calibrée sur le MTBF :
$$ p := \frac{2}{40\times365}, \qquad P_{\text{jour}} = \begin{pmatrix} 1-p & 0 & 0 \\ p & 1-p & 0 \\ 0 &
p & 1 \end{pmatrix}, \qquad p_{c,\cdot,\cdot,1,t} := P_{\text{jour}}^{60}. $$
Avec maintenance ($m=2$), $p_{c,x,x',2,t} = \mathbb{1}[x'=1]$ (retour certain à l'état neuf).

\paragraph{Coût de défaillance $\delta_{c,x,t}$.} Conformément à la Section 3.1 du document de référence,
$\delta_{c,x,t}$ est le coût attendu de $c$ dans l'état $x$ sur l'étape \emph{sans tenir compte des jours
gratuits}. Avec la pénalité de 100 par jour d'indisponibilité et sans météo, un composant en panne
($x=3$) laissé sans maintenance ($m=1$) immobilise le poste les 60 jours de l'étape :
$$ \delta_{c,x,t} = \begin{cases} 0 & x \in \{1,2\}, \\ 100 \times 60 = 6000 & x = 3. \end{cases} $$

\section{Colonnes et coût joint $\eta_{b,t}$}

Une colonne est un plan $b = (\text{comps}, \text{start\_days}, \text{free\_days})$ : l'ensemble des
composants maintenus, leur jour de démarrage dans l'étape, et le sous-ensemble de jours déclarés gratuits
($\kappa_b = |\text{free\_days}|$). Conformément à la Section 3.2 du document de référence, le coût joint
attribué à la colonne est \textbf{explicitement la somme des coûts sur les 60 jours de l'étape} :
$$ \eta_{b,t} = \sum_{\tau=1}^{60} p_\tau\,\big(\mathrm{am}_\tau - w_\tau\big) = 100\sum_{\tau=1}^{60}
\big(\mathrm{am}_\tau - w_\tau\big), $$
où $\mathrm{am}_\tau=1$ si au moins un composant de la colonne est en cours de maintenance le jour $\tau$,
et $w_\tau = \mathrm{am}_\tau \wedge f_\tau$ vaut 1 si ce jour est à la fois actif et déclaré gratuit (donc
sans pénalité). C'est bien $\eta_{b,t}$, et non $\delta$, qui porte le coût lié au planning
\emph{effectivement décidé} pour la période (jours de démarrage et jours gratuits choisis conjointement).

\section{Méthode 1 : problème maître restreint et génération de colonnes}

\subsection{Problème maître restreint (RMP)}

En notant $q_{c,x,m,k,\kappa,t}\ge0$ la probabilité jointe que $c$ soit dans l'état $x$, avec décision
$m\in\{1,2\}$, quota restant $k$ et jours gratuits utilisés $\kappa\le k$, et $B_{b,t}\ge0$ la probabilité
que la colonne $b$ soit choisie à l'étape $t$ :

\begin{align}
\min_{q,B}\quad & \sum_{c,x,(k,\kappa),t} \delta_{c,x,t}\, q_{c,x,1,k,\kappa,t}
 \ +\ \sum_{b,t} \eta_{b,t}\, B_{b,t}
 \ +\ \sum_{c,x,(k,\kappa),t} 100\, q_{c,x,2,k,\kappa,t} \label{eq:rmp-obj}\\
\text{s.c.}\quad
 & \sum_{m,\kappa\le k} q_{c,x,m,k,\kappa,t} \ge \sum_{x',m} p_{c,x',x,m,t-1} \!\!\sum_{\kappa'=0}^{10-k}
 q_{c,x',m,k+\kappa',\kappa',t-1}, && \forall c,x,k,\ t\ge2,\ t\not\equiv1\!\!\pmod6 \label{eq:rmp-flow1}\\
 & \sum_{m,\kappa\le10} q_{c,x,m,10,\kappa,t} \ge \sum_{x',m,(k',\kappa')} p_{c,x',x,m,t-1}\,
 q_{c,x',m,k',\kappa',t-1}, && \forall c,x,\ t\ge2,\ t\equiv1\!\!\pmod6 \label{eq:rmp-flow2}\\
 & q_{c,x,m,k,\kappa,t}=0, && \forall k\ne10,\ t\equiv1\!\!\pmod6 \label{eq:rmp-reset}\\
 & \sum_{m,\kappa\le10} q_{c,x,m,10,\kappa,1} = \alpha_{c,x}, && \forall c,x \label{eq:rmp-init}\\
 & \sum_{x,k\ge\kappa} q_{c,x,2,k,\kappa,t} \ \le\ \sum_{b:\,c\in C(b),\,\kappa_b=\kappa} B_{b,t},
 && \forall c,\kappa,t\quad [\lambda_{c,\kappa,t}] \label{eq:rmp-coupling}\\
 & q,B\ge0, \nonumber
\end{align}

où $\alpha_{c,x}$ est la distribution initiale (composants neufs, $x=1$), et la nouvelle année (toutes les
6 étapes) réinitialise le quota à $k=K_0=10$, conformément aux Sections 3.1--3.3 du document de référence.
La dernière contrainte est la contrainte de couplage entre l'occupation $q$ et les colonnes $B$ ; son dual
$\lambda_{c,\kappa,t}\le0$ sert au problème de pricing.

\subsection{Problème de pricing (MILP)}

Le coût réduit d'une colonne $b$ à l'étape $t$ est $\bar c(b,t) = \eta_{b,t} + \sum_{c\in C(b)}
\lambda_{c,\kappa_b,t}$. En introduisant, pour chaque composant $c$ : $z_c\in\{0,1\}$ (maintenu),
$m^{\text{start}}_{c,\tau}\in\{0,1\}$ ($\tau=1,\dots,60-d_c+1$, démarrage jour $\tau$) ; et globalement :
$f_\tau\in\{0,1\}$, $\mathrm{pay}_\tau\in[0,1]$, $\kappa^{\text{sel}}_\kappa\in\{0,1\}$ ($\kappa=0,\dots,10$,
exactement un actif), $y_{c,\kappa} = z_c \wedge \kappa^{\text{sel}}_\kappa$ (linéarisation de McCormick)
--- le MILP de pricing s'écrit :

\begin{align}
\min\quad & 100\sum_{\tau=1}^{60}\mathrm{pay}_\tau + \sum_{c=1}^{25}\sum_{\kappa=0}^{10}
\lambda_{c,\kappa,t}\, y_{c,\kappa} \label{eq:pricing-obj}\\
\text{s.c.}\quad
 & \textstyle\sum_{\tau=1}^{60-d_c+1} m^{\text{start}}_{c,\tau} = z_c, && \forall c=1,\dots,25
 \label{eq:pricing-c1}\\
 & \textstyle\sum_{c=1}^{25} \mathrm{ong}_{c,\tau} \ \le\ R=3, && \forall \tau=1,\dots,60
 \label{eq:pricing-c4}\\
 & \mathrm{pay}_\tau \ge \mathrm{ong}_{c,\tau} - f_\tau, && \forall c,\tau \label{eq:pricing-c5}\\
 & \textstyle\sum_{\kappa=0}^{10} \kappa^{\text{sel}}_\kappa = 1, \qquad \textstyle\sum_\tau f_\tau =
 \textstyle\sum_{\kappa=0}^{10} \kappa\,\kappa^{\text{sel}}_\kappa \label{eq:pricing-c6}\\
 & y_{c,\kappa} \le z_c,\ \ y_{c,\kappa}\le\kappa^{\text{sel}}_\kappa,\ \ y_{c,\kappa}\ge z_c+
 \kappa^{\text{sel}}_\kappa - 1, && \forall c,\kappa \label{eq:pricing-c7}
\end{align}
où $\displaystyle \mathrm{ong}_{c,\tau} := \!\!\sum_{\tau'=\max(1,\tau-d_c+1)}^{\min(\tau,60-d_c+1)}\!\!
m^{\text{start}}_{c,\tau'}$ est une expression linéaire de $m^{\text{start}}$, et non une variable
entière/compteur séparée. Ceci simplifie l'écriture du document de référence (qui introduit un compteur
$q^v_{c,\tau}$ de jours de travail restants, nécessaire quand la météo peut interrompre une opération en
cours) : sans aléa météo, une opération se déroule toujours exactement $d_c$ jours consécutifs sans
interruption, et puisque $\sum_\tau m^{\text{start}}_{c,\tau}=z_c\le1$ (au plus un jour de démarrage actif),
la formule ci-dessus donne automatiquement $\mathrm{ong}_{c,\tau}\in\{0,1\}$ sans qu'il soit nécessaire de
l'imposer par une contrainte supplémentaire.

\textbf{Simplification : fusion de $\mathrm{am}_\tau$ et $w_\tau$ en une seule variable $\mathrm{pay}_\tau$.}
Dans une première version, deux variables binaires par jour étaient utilisées : $\mathrm{am}_\tau$
(``au moins un composant en maintenance ce jour'', défini par $\mathrm{am}_\tau\ge\mathrm{ong}_{c,\tau}\
\forall c$) et $w_\tau$ (linéarisation de McCormick de $\mathrm{am}_\tau\wedge f_\tau$, 3 contraintes),
avec un objectif en $\mathrm{am}_\tau-w_\tau$. Ces deux variables peuvent être fusionnées en une seule,
$\mathrm{pay}_\tau$, valant 1 ssi un composant est en maintenance ce jour \emph{et} que le jour n'est pas
gratuit : c'est exactement \eqref{eq:pricing-c5}. Puisque $\mathrm{ong}_{c,\tau}$ et $f_\tau$ sont binaires,
$\max_c(\mathrm{ong}_{c,\tau}-f_\tau)\in\{-1,0,1\}$ ; comme $\mathrm{pay}_\tau$ est minimisé dans
l'objectif, il se fixe automatiquement à cette valeur (bornée à 0 par en dessous), sans qu'il soit
nécessaire de le déclarer binaire --- $\mathrm{pay}_\tau\in[0,1]$ continu suffit. Cette fusion supprime
$\mathrm{am}_\tau$, $w_\tau$ et les 3 contraintes de linéarisation de $w_\tau$ par jour, au profit d'une
seule contrainte \eqref{eq:pricing-c5} (même nombre que l'ancienne contrainte définissant
$\mathrm{am}_\tau$).

\textbf{Remarque.} Toutes les variables de décision sont binaires --- $z_c$, $m^{\text{start}}_{c,\tau}$,
$f_\tau$, $\kappa^{\text{sel}}_\kappa$, $y_{c,\kappa}$ --- \emph{à l'exception de $\mathrm{pay}_\tau$},
déclarée continue mais automatiquement entière à l'optimum pour la raison ci-dessus. Il ne s'agit donc pas
d'un MILP au sens général (variables mixtes entières/continues arbitraires), mais d'un programme presque
entièrement 0-1, avec une seule variable continue par jour dont l'intégralité découle directement de la
structure de l'objectif plutôt que d'être imposée : $\mathrm{ong}_{c,\tau}$ est une expression linéaire de
$m^{\text{start}}$ (et non une variable), et le compteur entier $q^v_{c,\tau}$ du document de référence
(nécessaire seulement en présence d'aléa météo) a été supprimé.

\eqref{eq:pricing-c4} est la contrainte de capacité des 3 techniciens (extension par rapport au document
de référence, qui n'en impose pas). \eqref{eq:pricing-c6} force à choisir exactement un nombre $\kappa$ de
jours gratuits, dans la limite du quota restant à cette étape ($\kappa^{\text{sel}}_\kappa=0$
automatiquement pour $\kappa>k_t$ puisque $\lambda_{c,\kappa,t}$ n'est défini/utile que pour $\kappa\le
k_t$). Ce MILP comporte, par composant, $\sim\!1+60=61$ variables binaires ($z_c$ et $m^{\text{start}}_{c,
\cdot}$), plus $11+25\times11+60\approx346$ variables globales (dont 60 continues, $\mathrm{pay}_\tau$),
soit $\approx 25\times61+346\approx 1\,900$ variables --- une taille modeste pour un solveur MILP générique.

\section{Méthode 2 : relaxation directe (contraintes en espérance)}

Conformément à la Section 5 du document de référence, on relâche l'intégralité du MILP de pricing
\eqref{eq:pricing-obj}--\eqref{eq:pricing-c7} en variables continues $z_{c,t},m^{\text{start}}_{c,\tau,t},
f_{\tau,t},\mathrm{pay}_{\tau,t}\in[0,1]$, une seule copie par étape $t$ (pas de distribution $B_{b,t}$ sur
des colonnes, pas de RMP ni de boucle de pricing) :

\begin{align}
\min_{q,z,\dots}\quad & \sum_{c,x,(k,\kappa),t} \delta_{c,x,t}\, q_{c,x,1,k,\kappa,t} + \sum_{c,t}
100\, z_{c,t} + \sum_{\tau,t} 100\,\mathrm{pay}_{\tau,t} \label{eq:relax-obj}\\
\text{s.c.}\quad & \text{contraintes de flot et d'initialisation de } \eqref{eq:rmp-flow1}\text{--}
\eqref{eq:rmp-init}, \nonumber\\
 & \sum_{x,k\ge\kappa} q_{c,x,2,k,\kappa,t} \ \le\ z_{c,t}, && \forall c,\kappa,t \label{eq:relax-c1}\\
 & \sum_{x,m,(k,\kappa)} \kappa\, q_{1,x,m,k,\kappa,t} = \sum_\tau f_{\tau,t}, && \forall t
 \label{eq:relax-c2}\\
 & \text{contraintes linéaires de planification } \eqref{eq:pricing-c1}\text{--}\eqref{eq:pricing-c5}
 \text{ (relâchées), indexées par } t. \nonumber
\end{align}

La contrainte de capacité des techniciens \eqref{eq:pricing-c4} devient $\sum_c r_c\,\mathrm{ong}_{c,\tau,t}
\le 3$ en variables continues : ce n'est plus l'existence d'un planning \emph{entier réalisable} qui est
exigée, mais seulement que le nombre \emph{attendu} de composants en maintenance un jour donné reste
$\le3$ --- d'où le nom ``contraintes en espérance''. Ce LP se résout directement, sans génération de
colonnes, au prix d'une borne potentiellement moins bonne (Remarque 3 du document de référence). Comme en
Méthode 1, $\mathrm{pay}_{\tau,t}$ remplace la paire $(\mathrm{am}_{\tau,t},w_{\tau,t})$ ; ici toutes les
variables étant déjà continues par construction, la fusion ne change rien à la nature du problème, mais
simplifie l'écriture et réduit le nombre de contraintes.

\section{Implémentation}

Les deux méthodes sont codées dans \texttt{test\_25\_composants.jl}, en reprenant la structure de
\texttt{script\_simple\_pole.jl} : à chaque étape, on résout un LP fluide sur un horizon roulant $H=6$
(un an), on ne retient que la décision de la première étape locale, on simule réellement l'étape (coût
joint exact, pas l'approximation $\delta/\eta$), on écrit le résultat dans les fichiers, et on recommence.
Les politiques déjà calculées pour un état $(x_0,k_0,t_0)$ rencontré sont mises en cache (en mémoire et
dans le fichier \texttt{policies\_25\_\{cg,relax\}.csv}) et relues sans être recalculées, comme le
\texttt{cache}/\texttt{collect\_policies} de \texttt{script\_simple\_pole.jl}. Le temps de calcul de
chaque nouvelle politique (en secondes) est mesuré et écrit en dernière colonne du fichier de politiques.

\subsection{Arrondi et heuristique de réparation (Méthode 2)}

La solution du LP relaxé à la première étape locale n'est en général pas entière. Elle est convertie en
un plan admissible de la façon suivante :
\begin{enumerate}
  \item \textbf{Maintenance} : $z_c$ est arrondi au plus proche ($z_c\ge0.5 \Rightarrow$ composant $c$
        maintenu).
  \item \textbf{Date de démarrage} : pour chaque composant retenu, le jour de démarrage est
        $\arg\max_\tau m^{\text{start}}_{c,\tau}$.
  \item \textbf{Jours gratuits} : les jours candidats sont ceux avec $f_\tau\ge0.5$ ; si leur nombre
        dépasse le quota $k_0$ réellement disponible à cette étape, on ne garde que les $k_0$ jours de
        plus forte valeur fractionnaire $f_\tau$ (réparation de la contrainte de quota).
  \item \textbf{Capacité technicien} : le planning des composants retenus (avec leurs dates de démarrage
        souhaitées) est ensuite réparé par un algorithme glouton de type \emph{list scheduling} : les
        composants sont traités par date de démarrage souhaitée croissante (les égalités étant départagées
        par confiance $z_c$ décroissante), et chacun est placé au plus tôt jour faisable
        $\tau\ge$ date souhaitée tel que les $R=3$ techniciens ne sont jamais simultanément dépassés sur
        toute la durée $d_c$ de son opération ; si aucun jour ne convient après la date souhaitée, les
        jours \emph{antérieurs} sont également essayés ; si aucun jour de la fenêtre $[1,60-d_c+1]$ ne
        convient, le composant est \textbf{reporté} (retiré du plan de cette étape --- il sera reconsidéré
        à l'étape suivante, sa dégradation ayant continué).
\end{enumerate}

\subsection{Pourquoi le polytope relâché n'est pas l'enveloppe convexe des actions jointes}

La relaxation $\bar{\mathcal{B}}$ utilisée en Méthode 2 est obtenue en relâchant \emph{seulement}
l'intégralité des variables qui définissent $\mathcal{B}_0$ (le système linéaire $Wb\le w$ lui-même est
conservé, cf. Section 5.1 du document de référence). Pour tout polytope combinatoire défini de cette
façon, on a toujours $\mathrm{conv}(\mathcal{B}_0)\subseteq\bar{\mathcal{B}}$ (toute action admissible, et
donc toute combinaison convexe d'actions admissibles, satisfait trivialement le système relâché), mais
l'inclusion réciproque n'a \emph{aucune raison} d'être vraie en général : $\bar{\mathcal{B}}$ peut contenir
des points fractionnaires qui ne sont la combinaison convexe d'\emph{aucune} collection d'actions
réellement admissibles de $\mathcal{B}_0$. L'égalité $\bar{\mathcal{B}}=\mathrm{conv}(\mathcal{B}_0)$
n'est garantie que si le système $Wb\le w$ est déjà ``idéal'' pour cet ensemble entier --- c'est
exactement la condition de la Proposition 4 du document \emph{Fluid Policies...} (``every vertex of
$\bar{\mathcal{B}}_0$ lies in $\mathcal{B}_0$''), qui est une condition suffisante mais nullement
automatique, et que rien ne garantit ici.

La raison structurelle du décalage est la suivante : $\mathcal{B}_0$ combine, pour chaque composant $c$,
une contrainte de type ``un seul choix à la fois'' (le simplexe $\Delta(\mathcal{A}_c)$, ici relâché en
$z_{c,t}\in[0,1]$ avec, par composant, un unique jour de démarrage et une unique valeur de $\kappa$ pour
toute l'étape), \emph{couplée} à une contrainte linéaire partagée entre composants (la capacité
technicien, ou ici la contrainte \eqref{eq:relax-c1} indexée par $\kappa$). Le relâchement traite chacune
de ces deux familles de contraintes indépendamment, alors que c'est précisément leur \emph{interaction
conjointe} qui borne le véritable ensemble des mélanges probabilistes réalisables : une combinaison
convexe d'actions entières doit correspondre à un \emph{même} tirage aléatoire partagé par tous les
composants (puisque $b=((a_m)_m,y,y')$ est une action \emph{jointe}, choisie une seule fois par étape),
alors qu'un point fractionnaire de $\bar{\mathcal{B}}$ peut satisfaire les contraintes en ``mélangeant''
indépendamment, coordonnée par coordonnée, des configurations qui ne proviendraient jamais d'un seul et
même tirage.

C'est exactement ce qui se produit dans l'exemple ci-dessous : \eqref{eq:relax-c1} borne la masse $m=2$
par la \emph{même} variable $z_{c,t}$ pour chaque $\kappa$ pris séparément. Un point réellement issu d'une
combinaison convexe d'actions entières qui maintiendrait toujours $c$ (chacune avec un $\kappa$ différent,
mais un unique $\kappa$ par action) donnerait $z_{c,t}=\sum_\kappa(\text{poids de l'action à ce }\kappa)=1$
: le poids total de maintenance doit se lire sur \emph{une seule} des contraintes-$\kappa$, celle
correspondant au $\kappa$ réellement tiré. Le LP relâché, lui, n'impose pas cette cohérence entre les
$\mathcal{K}+1$ contraintes : il les satisfait chacune séparément, avec la même variable $z_{c,t}=0.25$,
en répartissant fictivement de la masse ``maintenance'' sur les quatre valeurs de $\kappa$ à la fois ---
un comportement qu'aucune combinaison convexe de colonnes entières ($b\in\mathcal{B}_0$, une seule valeur
de $\kappa_b$ chacune) ne peut reproduire. C'est une illustration concrète, propre à cette instance, de
pourquoi $\bar{\mathcal{B}}\supsetneq\mathrm{conv}(\mathcal{B}_0)$ ici.

\subsection{Limite observée de la relaxation}

Un test direct sur une petite instance (un composant en panne, $x=3$, les autres neufs) illustre une
limite propre à la contrainte \eqref{eq:relax-c1} : celle-ci borne, \emph{pour chaque $\kappa$
séparément}, la masse $m=2$ par la même variable $z_{c,t}$ (exactement comme dans le document de
référence). Le LP peut donc répartir une probabilité de maintenance totale proche de 1 en petites parts
$z_{c,t}/(\mathcal{K}+1)$ sur chacune des $\mathcal{K}+1$ valeurs de $\kappa$, ce qui satisfait
\eqref{eq:relax-c1} avec un $z_{c,t}$ artificiellement petit (observé : $z_{c,t}=0.25$ alors que la
maintenance est en réalité clairement optimale, l'alternative coûtant $\delta=6000$ contre $\sim\!500$ en
la maintenant). L'arrondi au seuil 0.5 peut alors échouer à détecter un besoin de maintenance pourtant
réel. C'est une manifestation concrète de $Z^{\text{relax}}\le Z$ (Proposition 3 du document de
référence) : la relaxation peut sous-estimer significativement ce qui est réellement requis, et
l'heuristique de réparation ci-dessus ne corrige que la faisabilité du planning (capacité, quota), pas la
qualité de la décision d'arrondi elle-même.

\section{MILP ou énumération pour le problème de pricing ?}

\subsection{Ordre de grandeur de l'énumération}

Énumérer les sous-ensembles de composants maintenus $S\subseteq\{1,\dots,25\}$ représente $2^{25} =
33\,554\,432\approx3.35\times10^7$ possibilités. Pour chacune, il faudrait en plus résoudre le
sous-problème d'ordonnancement (placement des $d_c$, $c\in S$, sur 3 techniciens en parallèle sur 60 jours,
et choix du nombre de jours gratuits) pour évaluer $\eta_{b,t}$ et le coût réduit --- un sous-problème
lui-même combinatoire. Sur l'instance existante à 7 composants, $2^7=128$ : une énumération complète y est
instantanée.

\subsection{Pourquoi le MILP est nécessaire ici}

Avec $C=25$, en supposant (optimiste) $\sim\!10$--$30\,\mu s$ par sous-ensemble pour évaluer le
sous-problème d'ordonnancement associé, une énumération complète coûte
$$ 3.35\times10^7 \times (10\text{ à }30)\,\mu s\ \approx\ 5\text{ à }17\text{ minutes} $$
\emph{pour un seul appel} au pricing. Ce dernier doit être résolu à chacune des $\bar T=120$ étapes et à
chaque itération de génération de colonnes (quelques dizaines à centaines d'itérations) : le coût total de
l'énumération se chiffrerait en dizaines voire centaines d'heures, contre une fraction de seconde par
appel pour le MILP~\eqref{eq:pricing-obj}--\eqref{eq:pricing-c7} ($\approx1\,900$ variables, dont
$\sim\!1\,600$ binaires) résolu par un solveur exploitant la relaxation LP et le branch-and-cut.

\subsection{Conclusion}

Il fait donc sens de résoudre le problème de pricing par MILP plutôt que par énumération pour cette
instance à 25 composants : le nombre de sous-ensembles candidats ($\sim3.35\times10^7$), répété sur 120
étapes et de nombreuses itérations de génération de colonnes, rend l'énumération prohibitive, alors qu'un
MILP de taille modeste reste résolu quasi instantanément. Ce constat s'inverserait pour l'instance
existante à 7 composants ($2^7=128$), où l'énumération directe reste préférable en pratique.

\end{document}
